\section{Related Work}
\label{sec:related_work}

\textbf{Linear Weight Interpolation and Model Merging.}
Weight averaging was first popularized in federated learning under frameworks like Federated Averaging (FedAvg) \cite{fedavg, fedavg2}, where weights from models trained on decentralized partitions of the same dataset are averaged. Wortsman et al. \cite{wortsman2022model} extended this concept to fine-tuned models, introducing \emph{Model Soups} to show that averaging multiple fine-tuned models on a single task improves accuracy and robust generalization. Subsequently, Ainsworth et al. \cite{ainsworth2022git} and Stoica et al. \cite{stoica2023zipit} investigated permutation symmetries and alignment, developing methods like \emph{Git Re-Basin} to align weight coordinates before averaging. Ilharco et al. \cite{ilharco2022editing} generalized weight interpolation to multi-task learning by introducing \emph{Task Arithmetic} (TA). Under TA, a task vector is defined as the weight difference between a fine-tuned model and its pre-trained base model. Linear combinations of these task vectors can steer, add, or subtract specific model capabilities. Our work directly builds on the task vector paradigm, but seeks to resolve linear interference through spectral analysis.

\textbf{Resolving Parameter Interference.}
Multi-task model merging is frequently bottlenecked by \emph{destructive interference}, where conflicting parameter updates from different tasks cancel each other out, degrading downstream performance \cite{yu2023language}. To alleviate this, TIES-Merging \cite{yadav2023ties} introduces a three-step heuristic pipeline: (1) trimming task vectors based on parameter magnitude, (2) resolving sign conflicts via consensus voting, and (3) scaling the merged vectors. DARE \cite{yu2023dare} simplifies this heuristic by randomly dropping a high percentage of task vector parameters and scaling up the remaining ones, demonstrating that a significant portion of weight updates are redundant. While both TIES and DARE are training-free, they rely on heuristic, discrete, or randomized masking schemes. In contrast, SVS resolves interference by operating in the spectral domain of the weight updates, applying standard Singular Value Decomposition (SVD) and the Eckart-Young-Mirsky Theorem to prune redundant high-frequency dimensions.

\textbf{Optimization-Based Merging and Its Limitations.}
A parallel line of research attempts to resolve interference by formulating merging as an optimization problem. AdaMerging \cite{yang2023adamerging} seeks to find optimal layer-wise coefficients by minimizing the entropy of model predictions on validation streams or unlabeled test data at test time. FoldMerge (Neural Origami) \cite{foldmerge2025} introduces non-linear normalizing flow networks to bend and warp weight spaces into a shared, mergeable space. SyMerge \cite{symerge2025} leverages test-time adaptation of low-rank adapters to adjust merged representations dynamically.

While effective, these optimization-driven approaches introduce substantial test-time computational costs and auxiliary parameters, which can limit their applicability in resource-constrained environments. Furthermore, recent studies (e.g., SAIM \cite{saim2025}) demonstrate that optimizing complex layer-wise parameters on test streams can lead to overfitting to specific validation distributions, a phenomenon referred to as the ``Overfitting-Optimizer Paradox.'' In contrast, our proposed SVS method avoids these complex optimization pipelines, achieving competitive merging performance in closed form with negligible computational cost and without any trainable parameters.

\textbf{Data and Computation Constraints in Merging Settings.} To ensure a fair and scientifically rigorous comparison, we explicitly distinguish between \emph{unsupervised, data-free offline merging operators} (such as Task Arithmetic, TIES-Merging, DARE, and our proposed SVS) and \emph{test-time optimization-driven frameworks} (such as AdaMerging, FoldMerge, and SyMerge). Offline model merging operators operate under a strict zero-data and zero-training constraint: they merge task-specific experts post-hoc in a single closed-form step without accessing any labeled validation sets, downstream test streams, or unlabeled test-time data. In contrast, optimization-driven methods require active test-time streams of data, labeled development sets to search layer-wise coefficients, or several minutes of iterative GPU training (e.g., training FoldMerge's 2.6M normalizing flow parameters or SyMerge's adapters). Because accessing downstream test data or conducting test-time training violates the core premises of strict offline model merging, we restrict our empirical baseline comparison to the complete suite of training-free, data-free offline merging operators (Task Arithmetic, TIES-Merging, and DARE), which ensures a scientifically fair evaluation under identical resource and data limitations.

\textbf{Spectral Methods and Low-Rank Projections.}
Low-rank adaptations have become a standard approach for parameter-efficient fine-tuning (PEFT), exemplified by Low-Rank Adaptation (LoRA) \cite{hu2021lora}, which updates weights using low-rank product matrices during training. More recently, a nascent line of work has explored applying Singular Value Decomposition (SVD) post-hoc for model compression and merging. In particular, Gargiulo et al. \cite{gargiulo2025tsv} proposed \emph{Task Singular Vectors} (TSV), performing SVD on task vectors and introducing TSV-Compress to retain only the top singular components to denoise weight updates. Concurrently, Stoica et al. \cite{stoica2025svdmerging} explored SVD-based merging using whitening transformations to align weight spaces, while Ortho-Merge \cite{orthomerge2025} developed orthogonal subspace projection schemes to prevent multi-task interference. 

While SVS shares the post-hoc SVD-based projection formulation of TSV-Compress, our work is fundamentally distinguished by its focus and novel theoretical and algorithmic contributions:
\begin{enumerate}
    \item \textbf{Global Scaling Cancellation Theory:} Rather than treating SVD-based model merging as a purely empirical compression heuristic, we provide the first formal proofs of scale-invariance in normalized architectures (Section \ref{sec:scale_invariance_proof}), explaining mathematically why global scale-preservation algorithms are redundant.
    \item \textbf{Information-Theoretic Adaptation:} Unlike TSV or other SVD-based merging methods which apply uniform ranks across all layers, we propose \emph{Entropy-SVS} (Section \ref{sec:entropy_svs_methodology}), an information-theoretic scheme that dynamically allocates ranks across deep backbones using Shannon spectral entropy of singular values, tracing a superior Pareto frontier.
\end{enumerate}
Therefore, SVS serves as a simplified, purest-form offline baseline that isolates the core benefits of spectral-domain low-pass filtering while providing deep theoretical and adaptive algorithmic insights for the merging community.
